%% Clear sans, courier

% Klasa:
\documentclass[12pt]{mwart}

% Kodowanie:
\usepackage[utf8]{inputenc}

% Czcionki:
\usepackage[sfdefault]{ClearSans} %% option 'sfdefault' activates Clear Sans as the default text font
\usepackage[T1]{fontenc}
\usepackage{courier}
\usepackage[T1]{polski}

% Geometria i style
%\usepackage{showframe} % debug
\usepackage{geometry}
\geometry{a4paper,twoside=true,top=20mm,bottom=30mm,inner=35mm,outer=20mm}
\pagestyle{plain}
\usepackage{float}
\usepackage{afterpage}

% Grafika
\usepackage{graphicx}

% Dane:
\title{Uniwersalny inteligentny sterownik silników do zastosowań przemysłowych}
\author{Łukasz Przeniosło}

% Zmienne srodowiskowe:
\newcommand{\paperTitlePl}{UNIWERSALNY INTELIGENTNY STEROWNIK SILNIKÓW DO ZASTOSOWAŃ PRZEMYSŁOWYCH}
\newcommand{\paperTitleEn}{UNIVERSAL INTELIGENT MOTOR DRIVER FOR INDUSTRIAL APPLICATIONS}
\newcommand{\topicChooseDate}{01.08.2015 r.}
\newcommand{\topicAcceptationDate}{01.06.2016 r.}
\newcommand{\tocLessLeft}[1]{\textsc{\textbf{\large \flushleft{#1}}} \vspace{3mm} \\ \indent }
\newcommand{\tocLessCenter}[1]{\textsc{\textbf{\large \centerline{#1}}} \vspace{3mm} \\ \indent }

% Indeksowanie w pdfie wyjsciowym
\usepackage{hyperref}

% START:
\begin{document}
\thispagestyle{empty} % usun nr strony ze strony tytulowej

\begin{figure}[h]
	\centering
	\includegraphics[height=18mm]{grafiki/logo_zut.pdf}
	\hspace{5mm}
	\includegraphics[height=18mm]{grafiki/logo_zut_we.pdf} \\
	\vspace{20mm}


	{\bf \large Łukasz Przeniosło} \\ 
	\vspace{5mm}

	nr albumu: 21789 \\
	kierunek studiów: Elektrotechnika \\
	specjalność: urządzenia i instalacje elektryczne \\
	forma studiów: studia niestacjonarne \\
	\vspace{20mm}

	\textsc{\textbf{\large \paperTitlePl}} \\
	\vspace{5mm}

	{\bf \large \paperTitleEn} \\
	\vspace{45mm}
	
	
	Praca dyplomowa magisterska \\
	napisana pod kierunkiem: \\
	\vspace{2mm}
	
	{\bf \large dr inż. Michała Bonisławskiego} \\
	\vspace{2mm}
	
	Katedra Elektroenergetyki i Napędów Elektrycznych
	\vspace{5mm}
	
\end{figure}

\begin{flushleft}
	{\footnotesize Data wydania tematu pracy: \topicChooseDate} \\
	{\footnotesize Data złożenia tematu pracy: \topicAcceptationDate} 
\end{flushleft}

\vspace{\fill \centerline{Szczecin 2016}}
\clearpage

\afterpage{\null \thispagestyle{empty} \newpage}
\clearpage

\thispagestyle{empty}

\tocLessCenter{OŚWIADCZENIE AUTORA PRACY DYPLOMOWEJ} 

Oświadczam, że praca dyplomowa inżynierska/magisterska pt. \\
 „\paperTitlePl” \\
napisana pod kierunkiem: dr. inż. Michała Bonisławskiego \\
jest w całości moim samodzielnym autorskim opracowaniem sporządzonym przy wykorzystaniu wykazanej w pracy literatury przedmiotu i materiałów źródłowych.
Złożona w dziekanacie Wydziału Elektrycznego treść mojej pracy dyplomowej w formie elektronicznej jest zgodna z treścią w formie pisemnej. \\

Oświadczam ponadto, że złożona w dziekanacie praca dyplomowa ani jej fragmenty nie
były wcześniej przedmiotem procedur procesu dyplomowania związanych z uzyskaniem
tytułu zawodowego w uczelniach wyższych.
\vspace{30mm}


\begin{flushright}
	\begin{tabular}{c}
  		............................................\\
  		podpis dyplomanta
	\end{tabular}
\end{flushright}

\vspace{5mm}

\begin{flushleft}
	Szczecin, dn. \topicAcceptationDate{}
\end{flushleft}


\clearpage

\thispagestyle{empty}

\tocLessLeft{Streszczenie pracy}
Praca przedstawia projekt systemu sterowniczego do różnorodnych silników elektrycznych, na który składają się między innymi: Zbudowane w oparciu o 32 bitowy mikrokontroler urządzenie do którego podłączany jest silnik i peryferia (sterownik) oraz aplikacja PC-owa służąca do nastaw niezbędnych do poprawnego działania układu parametrów brzegowych oraz badania pracy zasilanego przez sterownik silnika. Rozwiązanie ma zostać zastosowane do sterowania elementami ruchomymi w maszynach typu {\it Pick and Place}, drukarkach szablonowych, drukarkach 3D oraz we wszelkiego rodzaju konwojerach. W części opisowej dokumentu, została zaprezentowana niezbędna teoria nt. silników elektrycznych oraz sterowania. Ponadto, zostały porównane komercyjne rozwiązania sterowników różnej klasy, które są dostępne na rynku.

\tocLessLeft{Słowa kluczowe}
Sterowanie, silnik, elektronika mocy, oprogramowanie

\tocLessLeft{Abstract}
The paper presents a project of a control system, that can be used with a various types of electric motors. Among others the main components of the system are: built upon a 32 bit microcontroller device, to which the motors and peripherals are connected (the driver) and a PC application used for defining critical operating parameters and for gathering connected motor work diagnostics . This application is to be used with Pick and Place machines, stencil printers, 3D printers and all sort of conveyors. In the generic part of the text, the essential theory of electric motors and control has been described. Also, commercially available on the market  motor drivers of various class has been compared.

\tocLessLeft{Key words}
Controll, motor, power electronics, software


\clearpage

\tableofcontents
\clearpage


\section*{Wprowadzenie}

Chcąc zaprojektować system napędowy dla określonej aplikacji należy wziąć pod uwagę bardzo wiele czynników. Najważniejsze z nich to: koszt i dostępność rozwiązań gotowych na rynku. Ceny rozwiązań przemysłowych bardzo rzadko mają odzwierciedlenie w zastosowanym w ich strukturze {\it hardwarze}, jest to raczej cena za {\it know-how} producenta, wiedzę ekspercką, wsparcie techniczne oraz dziesiątki roboczogodzin programistów i elektroników tworzących produkt. W przypadku posiadania dużych środków takie rozwiązanie wydaje się być najbardziej odpowiednie. Ma ono jednak jedną znaczącą wadę (pomijając cenę oczywiście). Pomimo tego że producenci sterowników do silników prześcigają się w ilości posiadanych cech, czy technologii, to korzystając z takiego urządzenia nijako jesteśmy zmuszeni z korzystania z zamkniętego, ograniczonego systemu. System taki może lepiej lub gorzej wpasowywać się w ramy naszej aplikacji. Rzadko kiedy wpasowuje się idealnie, a najczęściej brakuje mu jakichś właściwości, lub jest przewymiarowany. \

Drugą opcją jest zakup taniego zintegrowanego chipu, który po raz kolejny w zależności od wyłożonych środków, (w o wiele mniejszej skali) zawiera w sobie mniej lub więcej logiki i czasami końcówki mocy. Taki układ można także zbudować z elementów dyskretnych, co nie obniża znacząco kosztów, lecz daje możliwość większej kontroli nad parametrami układu. Wady takiego rozwiązania są takie, że zintegrowane sterowniki są projektowane pod zastosowanie z konkretnym rodzajem silnika (np tylko do silników krokowych) i służą do sterowania silnikami małej mocy jeśli mają zintegrowane końcówki. W przypadku stosowania układów tego typu oraz dyskretnych, projektant jest zmuszony zaprojektować obwód drukowany oraz napisać bardzo dużą ilość kodu zawierającego wiele skomplikowanych procedur, które muszą działać w czasie rzeczywistym. \

W przypadku, kiedy ilość produkowanych maszyn współpracującymi ze sterownikami do silników jest duża, posiadając odpowiednią wiedzę i umiejętności, bardziej opłaca się zaprojektować rozwiązanie własne, "od zera". Ponadto, proces tworzenia tego typu systemu dostarcza ogromną ilość wiedzy praktycznej.
  
\subsection*{Cel pracy}

Zbudować sterownik do sterowania różnego typu silników (krokowych, {\it BLDC}, liniowych i {\it DC}), z możliwością dostosowywania parametrów, używając aplikacji PC. Aplikacja ma pozwalać na strojenie działających w sterowniku regulatorów, wybierać tryby pracy i badać działanie systemu w czasie rzeczywistym. Urządzenie ma zostać użyte w maszynach przemysłowych służących do produkcji urządzeń elektronicznych.

\subsection*{Zakres pracy}

Przegląd i analiza rozwiązań technologicznych, układów elektronicznych i oprogramowania stosowanych w sterownikach silników oferowanych przez wybranych producentów. Projekt uniwersalnego, kompaktowego sterownika do silników {\it DC} (prądu stałego), {\it BLDC (Brushless DC Motor)}, krokowych oraz liniowych {\it VCM (Voice Coil Motor}, inaczej {\it VCA - Voice Coil Actuator}) wraz z opisem sposobów regulacji i algorytmów zawierających elementy {\it AI}, które zostaną zaimplementowane w nowoczesnym 32 bitowym mikrokontrolerze z rodziny {\it ARM}. Opracowanie aplikacji na PC umożliwiającej dostrajanie parametrów dla konkretnego silnika oraz analizowanie jego pracy. Wykonanie, uruchomienie i badania laboratoryjne prototypu urządzenia. Prezentacja przykładowego zastosowania sterownika działającego w rzeczywistym układzie.

\clearpage


\end{document}

